\chapter{Design}

This chapter outlines the architectural design decisions and data structures used in implementing the search algorithms within the Dataclenz library.

\section{System Architecture}

The search algorithm implementations are integrated into the Dataclenz library's broader data management framework. The library provides a DataFrame structure for data storage and manipulation, with search capabilities operating as methods on this structure.

The overall architecture follows a modular design pattern, separating concerns between:

\begin{itemize}
    \item Data storage and representation (DataFrame structure)
    \item Algorithm implementation (search functions)
    \item Performance measurement (metrics collection)
    \item Results analysis (visualization and reporting)
\end{itemize}

\section{Data Structures}

\subsection{DataFrame Structure}
The Dataclenz library uses a custom DataFrame implementation to store and manage data:

\begin{lstlisting}[style=c]
typedef struct {
    Column *columns;  // Array of columns
    char **column_names;  // Array of column names
    int num_columns;  // Number of columns
    int num_rows;  // Number of rows
    int max_rows;  // Maximum number of rows allocated
} DataFrame;
\end{lstlisting}

This structure provides a tabular representation of data with typed columns, similar to databases or spreadsheets.

\subsection{Column Structure}
Each column within a DataFrame contains data of a specific type:

\begin{lstlisting}[style=c]
typedef struct {
    void *data;  // Pointer to the column data
    ColumnType type;  // Data type of the column
    int length;  // Length of the column
} Column;
\end{lstlisting}

The Column structure uses a void pointer to accommodate different data types, with the ColumnType enumeration specifying the actual type:

\begin{lstlisting}[style=c]
typedef enum {
    TYPE_INT,
    TYPE_FLOAT,
    TYPE_STRING
} ColumnType;
\end{lstlisting}

\subsection{Performance Metrics Structure}
For measuring and recording algorithm performance, a dedicated structure was implemented:

\begin{lstlisting}[style=c]
typedef struct {
    char search_type[20];  // Name of search algorithm
    int dataset_size;  // Size of the dataset
    int num_searches;  // Number of searches performed
    double avg_search_time;  // Average search time
    int avg_comparisons;  // Average number of comparisons
    int successful_searches;  // Number of successful searches
    int memory_usage_kb;  // Memory usage in KB
} SearchPerformanceMetrics;
\end{lstlisting}

This structure encapsulates all relevant performance metrics for each test case, facilitating comprehensive analysis.

\section{Algorithm Flow}

\subsection{Binary Search Flow}
The Binary Search algorithm follows this general flow:

\begin{enumerate}
    \item Initialize search boundaries to cover the entire sorted array
    \item Calculate the middle point of the current search bounds
    \item Compare the target value with the middle element
    \item If equal, return the index of the middle element
    \item If target is less than the middle element, search the lower half
    \item If target is greater than the middle element, search the upper half
    \item Repeat steps 2-6 until target is found or search space is empty
\end{enumerate}

\subsection{Jump Search Flow}
The Jump Search algorithm follows this sequence:

\begin{enumerate}
    \item Calculate the optimal jump step size (typically sqrt(n))
    \item Jump ahead by step size until a value greater than the target is found
    \item Perform a linear search backward from the point where the jump ended
    \item Return the index if target is found, otherwise indicate not found
\end{enumerate}

\section{Flow Chart}

[THIS IS PLACEHOLDER FOR FLOW CHART - The flowchart would illustrate the step-by-step process of both search algorithms, including initialization, comparison steps, and termination conditions.]

\section{Integration with DataFrame Operations}

The search algorithms are integrated with the DataFrame structure through wrapper functions that:

\begin{enumerate}
    \item Verify the column exists and is of a compatible type
    \item Sort the data if required (both algorithms require sorted input)
    \item Call the appropriate search algorithm implementation
    \item Return indices or a filtered DataFrame containing matching rows
\end{enumerate}

This integration ensures the search operations work seamlessly with the rest of the Dataclenz library's functionality.

\section{Design Considerations}

Several key design considerations influenced the implementation:

\begin{itemize}
    \item \textbf{Generality:} The algorithms are designed to work with multiple data types
    \item \textbf{Performance:} Critical sections are optimized to minimize unnecessary operations
    \item \textbf{Memory Efficiency:} Both algorithms maintain O(1) space complexity
    \item \textbf{Error Handling:} Robust checking for edge cases and invalid inputs
    \item \textbf{Testability:} Structure allows for isolated testing of component parts
\end{itemize}

These considerations ensure the search implementations are both efficient and robust across a wide range of use cases.